\documentclass[10pt,a4paper]{amsart}
\usepackage[margin=19mm]{geometry}
\usepackage{amsmath,amssymb,mathtools}
\usepackage[scr=rsfs,cal=euler]{mathalfa}
\usepackage{xfrac}
\usepackage[unicode,hidelinks,bookmarks=false]{hyperref}
\newcommand{\ie}{i.e.\ }
\newcommand{\cf}{cf.\ }
\newcommand{\resp}{respectively\ }
\newcommand{\Subsection}{\subsection}
\providecommand{\nobreakdash}{\nobreak\hbox{-}}
\setlength{\emergencystretch}{2em}
\multlinegap0pt
\usepackage{mathtools}
\usepackage{amssymb}
\usepackage[shortlabels]{enumitem}
\newtheorem{lemma}{Lemma}
\newtheorem{coro}{Corollary}
\newtheorem{theorem}{Theorem}
\newcommand{\A}{\mathbb{A}}
\newcommand{\C}{\mathbb{C}} %is there a conflict somewhere?
\newcommand{\Ch}{\textup{Ch}}
\newcommand{\N}{\mathbb{N}}
\newcommand{\Spec}{\textup{Spec }}
\newcommand{\Z}{\mathbb{Z}}
\newcommand{\bm}{\mathbf{m}}
\newcommand{\cK}{\mathcal{K}}
\newcommand{\cM}{\mathcal{M}}
\newcommand{\cN}{\mathcal{N}}
\newcommand{\cO}{\mathscr{O}}
\newcommand{\cQ}{\mathcal{Q}}
\newcommand{\gr}{\textup{gr}}
\newcommand{\sA}{\mathscr{A}}
\newcommand{\shD}{\mathscr{D}}
\newcommand{\supp}{\textup{supp}}
\title{Generalized nearby cycles via relative and logarithmic $\mathscr{D}$-modules}
\author{Lei Wu}
\date{Source: arXiv:2602.05314v3 (CC BY 4.0)}
\begin{document}
\maketitle

\noindent\emph{Source and licence.} Generalized nearby cycles via relative and logarithmic $\mathscr{D}$-modules. Version arXiv:2602.05314v3 (CC BY 4.0), distributed by arXiv under the Creative Commons Attribution 4.0 International licence, http://creativecommons.org/licenses/by/4.0/, as recorded in the arXiv licence metadata for this exact version and shown on the abstract page https://arxiv.org/abs/2602.05314v3. Full licence notice: \texttt{licenses/arxiv-2602.05314-CC-BY-4.0.txt}.\par\smallskip

\noindent\textit{Editorial scope.} Counted unit: the article's theorem thm:nozeroproperBSI (source0.tex lines 5713--5727). The article's complete original proof of it is delivered with it (source0.tex lines 5729--5938, one \texttt{proof} environment of 210 lines). No mathematics is written, repaired or re-derived: the statement and the proof are the article's own lines, in the article's own notation, and no retained line is substituted or re-flowed.  Retained uncounted as the source-local context the proof needs: lines 519--521; lines 679--682; lines 692--730; lines 2144--2155; lines 2354--2393; lines 4428--4439; lines 5619--5630.  Nothing was omitted from any retained span, so the package carries no omission record.  No retained line contains a citation, so the wrapper carries no bibliography and no bibliography entry.  No retained line contains a figure environment, an image-inclusion command or drawing code, so no image is delivered and the package declares no asset dependency.  Editorial formatting only, in addition: an AMS \texttt{amsart} preamble that restates the article's own declarations and macros used by the retained lines, namely the environments \texttt{lemma}, \texttt{coro}, \texttt{theorem} on separate counters, together with the standard packages those lines need.  The retained environments print in this document as Lemma 1, Coro 1, Theorem 1 in order of appearance, whereas the article prints the same items with the numbering of its own full document.  The complete native source of the article as distributed in the arXiv e-print for this exact version is bundled unchanged as \texttt{sources/194026/source0.tex}.
\begin{equation}\label{eq:uncanincl}
    \sA_X^r\hookrightarrow \sA_{Y,D}^k\hookrightarrow \shD_{Z,E},
\end{equation}

\begin{lemma}\label{lm:twosidedideal}
If $K\subseteq \N^r$ is a monoid ideal, then
$\shD_{Z,E}\cdot \alpha(K)\subseteq \shD_{Z,E}$ is a two-sided ideal.
\end{lemma}

\begin{lemma}\label{lm:logsp}
We have a natural isomorphism of $\C$-algebras
\[
i_Y^*(\shD_{Z,E})
=
\frac{\shD_{Z,E}}{\shD_{Z,E}\cdot(t_{l+1},\dots,t_r)}
\simeq
\sA^k_{Y,D},
\]
such that the inclusion $\sA^k_{Y,D}\hookrightarrow \shD_{Z,E}$ in \eqref{eq:uncanincl} is a section of the natural $\C$-algebra morphism
\[
\shD_{Z,E}\longrightarrow i_Y^*(\shD_{Z,E}).
\]
Moreover, we have the induced filtration on $\sA^k_{Y,D}$:
\[
F^\sharp_\bullet\sA^k_{Y,D}
\simeq
i_Y^*(F^{\log}_\bullet\shD_{Z,E}),
\]
so that
\[
\gr^\sharp_\bullet\sA^k_{Y,D}
\simeq
\gr^{\log}_\bullet\shD_{Y,D}\otimes_\C \C[s_{l+1},\dots,s_r].
\]
In particular, when $k=r$, we have
\[
i_X^*(\shD_{Z,E})\simeq \sA_X^r
\quad \textup{and} \quad
F^\sharp_p\sA_X^r
=
\sum_{q+|\mathbf{i}|=p}F_q\shD_X\cdot \mathbf{s}^{\mathbf{i}},
\]
where
\[
i_X\colon X\hookrightarrow Z,\quad x\mapsto (x,0,\dots,0)
\]
is the closed embedding.
\end{lemma}

\begin{coro}\label{cor:puregradequot}
If $\cN$ is an $\sA_{Y,D}^R$-lattice such that $\cN(*D)$ is a relative regular holonomic $\sA_Y^R$-module, then
\[
\Ch^\dagger(\cN/\alpha(\bm)\cdot\cN)
=
\Ch^\dagger(\cN)|_{(\alpha(\bm)=0)}.
\]
In particular, if the support of $\cN$, as an $\cO_Y$-module, intersects $(\alpha(\bm)=0)$ nontrivially for $\bm\in \N^l$, then
\[
j(\cN/\alpha(\bm)\cdot\cN)=j(\cN)+1.
\]
\end{coro}

\begin{theorem}\label{thm:mainsuppfiberation}
Let $\cN$ be a $j$-pure $\sA^R_{Y,D}$-lattice of GE-type with respect to
$G=(g_1,\dots,g_l)$
such that $\cN(*D)$ is a relative regular holonomic $\sA^R_Y$-module. Then $\cN$ is a relative holonomic $\sA_X^{R\otimes_\C R_l}$-module, and $\cN(*D)$ is a relative regular holonomic $\sA_X^{R\otimes_\C R_l}$-module.

Moreover, for nonzero $\bm\in\N^l$, the projection
\[
p_{\log}\colon \Spec(R\otimes_\C R_l)\to \Spec R
\]
induces a surjective morphism $\bar p_{\log}
=
p_{\log}|_{\supp^h_X(\cN/\alpha(\bm)\cdot\cN)}$,
\[
\bar p_{\log}
\colon
\supp^h_X(\cN/\alpha(\bm)\cdot\cN)
\longrightarrow
\supp_Y(\cN(*D))
=
\supp^R_{Y\setminus D}(\cN|_{Y\setminus D})
\]
such that, for every closed point
$a\in \supp_Y(\cN(*D)),$
the fiber
\[
\bar p_{\log}^{-1}(a)\subseteq \Spec R_l
\]
is a finite union of codimension-one translated linear subvarieties of the form
\[
\left(\sum_{i=1}^l a_i s_i+b=0\right),
\]
where $(a_1,\dots,a_l)\in \Z_{\ge 0}^l$ is a nonzero primitive vector satisfying
$\bm\cdot(a_1,\dots,a_l)\neq 0.$
Here
$\supp^h_X(\cN/\alpha(\bm)\cdot\cN)$
denotes the top-dimensional part of
\[
\supp_X^{R\otimes_\C R_l}(\cN/\alpha(\bm)\cdot\cN).
\]
\end{theorem}

\begin{lemma}\label{lm:logKashequiv}
Let $\cM$ be a finitely generated $\sA^R_{Y,D}$-module supported on
$D^I=Y_{I^c}$
for some $I\subseteq \{1,\dots,l\}$. Then $\cM$ is a finitely generated
$\sA_{Y_{I^c},\bar D_{I^c}}^{R\otimes_\C R_I}$-module via the inclusion
\[
\sA_{Y_{I^c},\bar D_{I^c}}^{R\otimes_\C R_I}
\hookrightarrow
\sA^R_{Y,D}
\]
induced by \eqref{eq:uncanincl}.
\end{lemma}

\begin{lemma}\label{lm:subbahgb}
If $\cN(*E)$ is a holonomic $\shD_Z$-module, then $B^K(\cN)$ contains a
nonzero polynomial $b=b(s_1,\dots,s_r)$ of the form
\[
b
=
\prod_{(L,\alpha)}
\big(L\cdot(s_1,\dots,s_r)+\alpha\big)^{n_{L,\alpha}},
\]
where $L\in \Z_{\ge 0}^r$, $\alpha\in \C$, and only finitely many exponents
$n_{L,\alpha}$ are nonzero.
\end{lemma}

\begin{theorem}\label{thm:nozeroproperBSI}
Let
$F=(f_1,\dots,f_r)\colon X\to \A_\C^r$
be a regular map, and let $K\subsetneq \N^r$ be a proper monoid ideal. Suppose
that $\cN(*E)$ is a nonzero regular holonomic $\shD_Z$-module supported on the
graph of $F$. If the support of $\cN$, as an $\cO_Z$-module, intersects
$Z_K$ nontrivially, then $B^K(\cN)$ is a nonzero proper ideal. In particular,
if
$(f_1=f_2=\cdots=f_r=0)\subseteq X$
is nonempty, then
\[
B_F^K\subseteq \C[s_1,\dots,s_r]
\]
is a nonzero proper ideal.
\end{theorem}

\begin{proof}
By Lemma \ref{lm:subbahgb}, it remains to prove that
$\cN|_{Z_K}\neq 0.$
Since the support of $\cN$ is a Zariski closed subset of $Z$, we may assume
that the support of $\cN$ meets $Z_{K_I}$ for some nonempty subset
$I\subseteq \{1,\dots,r\}$
such that
$Z_{K_I}\hookrightarrow Z_K.$
It is therefore enough to prove
$\cN|_{Z_{K_I}}\neq 0$
under the assumption that the support of $\cN$ intersects $Z_{K_I}$
nontrivially.

For simplicity, assume
$I=\{1,\dots,k\}$
for some $k\le r$, and write
\[
I_j=\{j,j+1,\dots,r\}
\]
for $j\le k$. We argue by induction on $|I|$ to prove that
$\cN|_{Z_{K_I}}\neq 0.$
By the inductive hypothesis, we may assume that
\[
\cN_j
\coloneqq
\cN|_{Z_{K_{\bar I_j}}}
\neq 0
\]
for $j<k$, and that the support of $\cN_j$, as an $\cO_Z$-module or as an
$\cO_{Y_j}$-module, intersects $Z_{K_I}$ nontrivially.

By Lemma \ref{lm:logsp}, $\cN_j$ is a finitely generated
$\sA_{Y_j,\bar E_j}^{R_{j-1}}\text{-module},$
where
\[
Y_j=Y_{\bar I_j}
\quad \text{and} \quad
\bar E_j=E_{I_j}\cap Y_j.
\]
By Theorem \ref{thm:mainsuppfiberation}, applied to the log structure
\[
\N^{j-1}
\hookrightarrow
S[t_j^{\pm1},\dots,t_r^{\pm1}][\N^{j-1}],
\]
the module $\cN_j(*\bar E_j)$ is relative regular holonomic over
$\sA_{Y_j}^{R_{j-1}}$.

Consider the short exact sequence
\[
0
\to
\cK_j
\to
\cN_j
\to
\cQ_j
\to
0,
\]
where $\cK_j$ and $\cQ_j$ are, respectively, the kernel and image of the
natural map
\[
\cN_j\to \cN_j(*\bar E_j)=\cN_j(*E_{I_j}).
\]
Then $\cQ_j$ is a lattice, and hence has no $t_i$-torsion for $i\ge j$.
Therefore we have a short exact sequence
\[
0
\to
\frac{\cK_j}{t_j\cK_j}
\to
\cN_{j+1}
\to
\frac{\cQ_j}{t_j\cQ_j}
\to
0.
\]

If the $\cO_{Y_j}$-support of $\cQ_j$ intersects $Z_{K_I}$ nontrivially, then
Corollary \ref{cor:puregradequot} gives
\[
\Ch^\dagger\left(\frac{\cQ_j}{t_j\cQ_j}\right)
=
\Ch^\dagger(\cQ_j)|_{(t_j=0)}.
\]
In particular, $\cN_{j+1}\neq 0$, and the support of $\cN_{j+1}$ intersects
$Z_{K_I}$ nontrivially.

Now suppose that the $\cO_{Y_j}$-support of $\cQ_j$ does not intersect
$Z_{K_I}$. Then the $\cO_{Y_j}$-support of $\cK_j$ must intersect $Z_{K_I}$
nontrivially by the inductive hypothesis. Consider the natural map
$\cK_j\to \cK_j(*D_{\widehat{j+1}})$,
where
\[
D_{\widehat{j+1}}
=
\left(
\prod_{\substack{i\in\{j,j+1,\dots,r\}\\ i\neq j+1}}t_i=0
\right)
\subseteq Y_j.
\]
Let $\cQ_j^1$ denote its image and $\cK_j^1$ its kernel. Then we have a short
exact sequence
\begin{equation}\label{eq:aaabbbses123}
0
\to
\cK_j^1
\to
\cK_j
\to
\cQ_j^1
\to
0.
\end{equation}
By construction, $\cQ_j^1$ is supported on
\[
D_{j+1}=(t_{j+1}=0)\subseteq Y_j,
\]
and $\cK_j^1$ is supported on $D_{\widehat{j+1}}$.

For $l=1,\dots,r-j$, set
\[
Y_j^l\coloneqq D_{j+l}
\quad \text{and} \quad
\bar D_{\widehat l}\coloneqq D_{\widehat{j+l}}\cap Y_j^l.
\]
By Lemma \ref{lm:logKashequiv}, $\cQ_j^1$ is an
$\sA_{Y_j^1,\bar D_{\widehat 1}}^{R_{j-1}[s_{j+1}]}\text{-lattice}.$
Since $\cN$ is relative regular holonomic over $\sA_X^{R_r}$ by
Theorem \ref{thm:mainsuppfiberation}, and since $\cQ_j^1$ is supported on the
graph of
\[
(f_j,f_{j+2},f_{j+3},\dots,f_r)
\]
as an $\sA_{Y_j^1,\bar D_{\widehat 1}}^{R_{j-1}[s_{j+1}]}$-module, Kashiwara
equivalence implies that
$\cQ_j^1(*\bar D_{\widehat 1})$
is relative regular holonomic over
$\sA_{Y_j^1}^{R_{j-1}[s_{j+1}]}.$

If the $\cO$-support of $\cQ_j^1$ intersects $Z_{K_I}$ nontrivially, then the
same argument applied to $\cQ_j$ shows that
${\cK_j}/{t_j\cK_j}\neq 0$,
and that the support of $\cN_{j+1}$ intersects $Z_{K_I}$ nontrivially.

If instead the $\cO$-support of $\cQ_j^1$ does not intersect $Z_{K_I}$, then
the $\cO$-support of $\cK_j^1$ intersects $Z_{K_I}$ nontrivially. Repeating
the same construction with
\[
\cK_j^1\to \cK_j^1(*D_{\widehat{j+2}}),
\]
we obtain another short exact sequence
\[
0
\to
\cK_j^2
\to
\cK_j^1
\to
\cQ_j^2
\to
0,
\]
where $\cQ_j^2$ is an
$\sA_{Y_j^2,\bar D_{\widehat 2}}^{R_{j-1}[s_{j+2}]}\text{-lattice}$
and
$\cQ_j^2(*\bar D_{\widehat 2})$
is relative regular holonomic over
$\sA_{Y_j^2}^{R_{j-1}[s_{j+2}]}$.
If the $\cO$-support of $\cQ_j^2$ intersects $Z_{K_I}$ nontrivially, then
${\cK_j^1}/{t_j\cK_j^1}\neq 0$,
and the support of this quotient intersects $Z_{K_I}$ nontrivially. Hence the
same is true for $\cK_j/t_j\cK_j$, and therefore for $\cN_{j+1}$.

Otherwise, we continue this process with
\[
\cK_j^2,\cK_j^3,\dots,\cK_j^{r-j}
\]
as necessary. In the worst case, the $\cO$-support of $\cQ_j^{r-j}$ does not
intersect $Z_{K_I}$. Then the $\cO$-support of $\cK_j^{r-j}$ intersects
$Z_{K_I}$ nontrivially. By construction, $\cK_j^{r-j}$ is supported on
$D_j$. Therefore, by Lemma \ref{lm:twosidedideal}, there exists $m\gg 0$ such
that
$t_j^m\cdot \cK_j^{r-j}=0$.
Working algebraically locally at a point in the intersection of the
$\cO$-support of $\cK_j^{r-j}$ with $Z_{K_I}$, it follows that
\[
\frac{\cK_j^{r-j}}{t_j\cK_j^{r-j}}\neq 0
\]
and that the support of this quotient intersects $Z_{K_I}$ nontrivially.
Since each $\cQ_j^a$ is a lattice, we have successive injections
\[
\frac{\cK_j^{r-j}}{t_j\cK_j^{r-j}}
\hookrightarrow
\frac{\cK_j^{r-j-1}}{t_j\cK_j^{r-j-1}}
\hookrightarrow
\cdots
\hookrightarrow
\frac{\cK_j}{t_j\cK_j}.
\]
Hence the supports of both
\[
\frac{\cK_j}{t_j\cK_j}
\quad \text{and} \quad
\cN_{j+1}
\]
intersect $Z_{K_I}$ nontrivially in this worst case as well. This completes
the induction and hence the proof.
\end{proof}

\end{document}
